\section{Lifting of Gaussian and Jacobian Sum}
So using the discussion on \emph{Dirichlet $L$-functions} we can relate Gaussian sums in $\bbF_q$ with Gaussian sum of lifted characters to any finite extension of $\bbF_q$. 
\begin{Theorem}{Davenport-Hasse Theorem}{thm:davenport-hasse}
  Let $\chi\in\Xq$ be an additive character and $\psi\in\Mq$ be a multiplicative character of $\bbF_q$, not both of them trivial. Suppose $\chi$ and $\psi$ are lifted to character $\chi^{(r)}$ and $\psi^{(r)}$ respectively of the finite extension field $E$ of $\bbF_q$ with $[E\colon \bbF_q]=r$. Then $$G(\psi^{(r)},\chi^{(r)})=(-1)^{r-1}G(\psi,\chi)^r$$
\end{Theorem}
\begin{proof}
  Like previous section $G$ is the set of all rational functions $r(X)=\nicefrac{h_1}{h_2}$ where $h_1,h_2\in\Phi$, $h_2$ non-zero. Let $\ovG$ be the subgroup of $G$ consisting of rational functions $r(X)=\frac{h_1(X)}{h_2(X)}$ having $h_1(0),h_2(0)\neq 0$. Then define the multiplicative function $\lm:G\to \bbC$ such that for all $r\in G$, $|\lm(r)|\leq 1$ and $\lm(1)=1$ in the following way: 

  If $r(X)\in G$ where $r(X)=\frac{h_1(X)}{h_2(X)}$ and in the common splitting field $E$ of $h_1,h_2$, they factor as: $$h_1(X)=\prod_{i=1}^{d_1}(X-\alpha_i)\qquad \text{and}\qquad h_1(X)=\prod_{j=1}^{d_2}(X-\beta_j)$$ where $\alpha_i,\beta_j\in E$ for all $i\in[d_1]$, $j\in[d_2]$ then $$\lm(r)=\psi\lt(\prod_{i=1}^{d_1}\alpha_i \cdot \prod_{j=1}^{d_2}\beta_j^{-1}\rt)\cdot \chi\lt(\sum_{i=1}^{d_1}\alpha_i-\sum_{j=1}^{d_2}\beta_j\rt)$$Now for any monic polynomial $h\in\Phi$ let $h(X)=X^n-c_1\cdot X^{n-1}+\cdots +(-1)^nc_n$ where $c_i\in\bbF_q$ for all $i\in[n]$. Then we have $\lm(h)=\psi(c_n)\cdot \chi(c_1)$. Clearly $\lm$ is a multiplicative map on $G$. Now for a given $(c_1,c_k)$ pair there are $q^{n-2}$ polynomials where same $(c_1,c_n)$ occurs. Hence $$\sum_{h\in\Phi_k}\lm(h)=q^{n-2}\sum_{c_1\in\bbF_q}\sum_{c_n\in\bbF_q^*}\psi(c_n)\cdot \chi(c_1)=q^{n-2}\lt(\sum_{c_n\in\bbF_q}\rt)\lt(\sum_{c_1\in\bbF_q}\chi(c_1)\rt)$$Since $\psi$ and $\chi$ are non-trivial by \thmref{thm:add-char-props} and \thmref{thm:mult-char-props} we have $\sum_{h\in\Phi_k}\lm(h)=0$ for $n>1$. 

  Now $\Phi(0)=\{1\}$ and $\Phi(1)=\{x-\alpha\colon \alpha\in\bbF_q^*\}$. Hence $$\sum_{h\in\Phi_1}\lm(h)=\sum_{\alpha\in\bbF_q^*}\psi(\alpha)\cdot \chi(\alpha)=G(\psi,\chi)$$Therefore $$\ovL(z,\lm)=1+G(\psi,\chi)\cdot z\qquad\text{and}\qquad L(s,\chi)=1+q^{-s}\cdot G(\psi,\chi)$$Therefore root of $\ovL(z,\lm)$ is $-G(\psi,\chi)$.  

  Now we will calculate $L_r$  where $r=[E\colon \bbF_q]$. From the definition of $L_r$ in previous section we have $$L_r =  \sum_{\substack{h\in\irr(\Phi),\\ \deg(h)\mid r}}\deg(h)\cdot \chi(h)^{\nicefrac{r}{\deg(h)}}= \sum_{\substack{h\in\irr(\Phi),\\ \deg(h)\mid r}} \deg(h)\cdot \lm\lt( h^{\nicefrac{r}{\deg(h)}}\rt)$$Now for any irreducible $h$ let $\gm$ be a root of $h$ in the splitting field of $f$. Let $$h^{\nicefrac{r}{\deg(f)}}= x^r-c_1\cdot x^{r-1}+\cdots+(-1)^k c_r=\prod_{i=o}^{r-1}\lt(X-\gm^{q^i}\rt)$$Hence 
  $c_1=\tr_{E/\bbF_q}(\gm)$ and $c_r=\norm_{E/\bbF_q}(\gm)$. Hence we have $$\lm\lt(h^{\nicefrac{r}{\deg(h)}}\rt)=\psi(c_r)\cdot \chi(c_1)=\psi\lt(\norm_{E/\bbF_q}(\gm)\rt)\cdot \chi\lt(\tr_{E/\bbF_q}(\gm)\rt)=\psi^{(r)}(\gm)\cdot \chi^{(r)}(\gm)$$Now we can express $\deg(h)\cdot \lm\lt(h^{\nicefrac{r}{\deg(h)}}\rt)=\sum_{\gm\in E,h(\gm)=0}\psi^{(r)}(\gm)\cdot \chi^{(r)}(\gm)$. 
  Therefore $$L_r= \sum_{\substack{h\in\irr(\Phi),\\ \deg(h)\mid k}} \deg(h)\cdot \lm\lt( h^{\nicefrac{k}{\deg(h)}}\rt)=\sum_{\substack{h\in\irr(\Phi),\\ \deg(h)\mid k}}\sum_{\gm\in E,h(\gm)=0}\psi^{(r)}(\gm)\cdot \chi^{(r)}(\gm) $$Now all irreducible polynomials of $\bbF_q$ whose degree divides $r$ with non-zero constant their roots in $E$ are disjoint sets and partition $E^*$. Therefore we have $$L_r=\sum_{\gm\in E^*}\psi^{(r)}(\gm)\cdot \chi^{(r)}(\gm)=G(\psi^{(r)},\chi^{(r)})$$. Since  by \thmref{thm:l-function-root-factorisation} we have $L_r=-(-G(\psi,\chi))^r$. So we have the theorem. 
\end{proof}